%MẪU CUỐN SÁCH TRONG VieTeX
%Dùng Vietex 2.1. với phông Unicode
%Người soạn : Nguyễn Hữu Điển, ĐHKHTN, ĐHQG HN
%Mail: huudien@vnu.edu.vn, CQ: (84 - 4) 557 2869
%NR: (84 - 4) 641 8848, DĐ: 0989061951
%%%%%%%%%%%%%%%%%%%%%%%%%

%-[Phần mở đầu chương
\chapter*{Những kí hiệu}
\markboth{{\it Phương pháp đại lượng bất biến}}{{\it Kí hiệu}}
\addcontentsline{toc}{section}{{\bf Những kí hiệu}}

Trong cuốn sách này ta dùng những kí hiệu với các ý nghĩa xác định trong bảng dưới đây:

\begin{center}
\begin{tabular}{ l  l }
$\mathbb N$ & tập hợp số tự nhiên\\
$\mathbb N^*$ & tập hợp số tự nhiên khác $0$\\
$\mathbb Z$ & tập hợp số nguyên\\
$\mathbb Q$ & tập hợp số hữu tỉ\\
$\mathbb R$ & tập hợp số thực\\
$\mathbb C$ & tập hợp số phức\\
$\equiv$& dấu đồng dư\\
$\infty$& dương vô cùng (tương đương với $+\infty$)\\
$-\infty$& âm vô cùng\\
$\emptyset$& tập hợp rỗng\\
$\binom{m}{k}$ & tổ hợp chập $k$ của $m$ phần tử\\
$\vdots$& phép chia hết\\
$\kchia$& không chia hết\\
$UCLN$& ước số chung lớn nhất\\
$BCNN$& bội số chung nhỏ nhất\\
$\deg$& bậc của đa thức\\
%MO&Mathematics Olympiad \\
IMO&International Mathematics Olympiad \\
APMO&Asian Pacific Mathematics Olympiad\\
%USAMO&International Mathematics Olympiad \\
\end{tabular}
\end{center}
